\section{Operational Implications}

The released artifact is a research prototype rather than a deployment guide, but the current results already expose several system-level constraints that any practical use would have to respect. Making these implications explicit helps keep the paper honest about what Ghostext does well and what would still need engineering beyond the present code release.

\paragraph{Shared state is not a minor setup detail.} Ghostext requires Alice and Bob to share more than a passphrase. They need the same prompt, the same backend family, the same tokenizer behavior, the same seed, and the same candidate-policy and codec settings. In the artifact we treat this as a virtue because it makes replay auditable. In a real communication workflow, however, it means that state distribution and upgrade management become part of the covert-channel problem. A silent model update or tokenizer drift is not just a performance nuisance; it can break communication altogether. Any future deployment-oriented design would therefore need an explicit strategy for version pinning, prompt management, and controlled migration between compatible configurations.

\paragraph{Payload size should be treated as an application parameter.} The evaluation shows that message length affects runtime and packet size much more clearly than it affects average bits/token. Short messages are cheaper to encode and decode, while long messages push both total token count and retry-driven tail latency upward. This means Ghostext is better viewed as a medium-throughput covert transport for short or moderate payloads than as a bulk hidden channel. If a future system needed stronger traffic uniformity, it would likely need message chunking, padding, batching, or a higher-level scheduling policy that decides when to send one packet versus several smaller packets. None of those channel-management mechanisms are in scope for the current artifact, but the released numbers make clear why they would matter.

\paragraph{Retry policy is a real throughput knob.} The encoder's retry budget is not only a safety net. It is part of the throughput/correctness trade-off. A stricter policy that stops early would reduce latency variance but would fail more often in low-entropy or unstable-candidate contexts. A more permissive policy can preserve recoverability at the cost of longer and less predictable encode times. Our released settings prefer successful authenticated recovery over tight latency control, which is why the 10x10 grid contains a visible long tail without any decode failures. This is a defensible choice for a correctness-first artifact, but it also means that future work should treat retry policy as a first-class systems parameter rather than as an implementation footnote.

\paragraph{Artifact release and deployment pressure point in opposite directions.} For reproducibility, we intentionally release prompt sets, message sets, full cover texts, hashes, runtime metadata, and exact scripts. That is the right choice for a paper artifact. A real covert deployment would often want the opposite: prompt diversity, less repeated structure, and fewer repeated public templates. This tension does not invalidate the artifact; it simply means that our released evaluation should be read as a transparent lower-level protocol study rather than as a deployment-ready operational recipe. The paper's contribution is that it makes the system measurable and auditable first, so that stronger deployment questions can be asked on top of a concrete baseline instead of on top of hand-wavy intuition.
